\noindent\textbf{Evidence note.} These are manual historical trace slices; `success` below means only the saved attempt record's compiler-verification outcome. No API provenance is claimed.
\paragraph{arith\_001 (manual trace; attempts 1 and 2 both rejected).}
\begin{quote}\small
\textbf{Initial candidate excerpt:}\par\ttfamily
import algebra.group\_power\\
import algebra.ring\\
\\
theorem add\_sq\_nat (a b : ℕ) : (a + b) \textasciicircum{} 2 = a \textasciicircum{} 2 + 2 * a * b + b \textasciicircum{} 2 := by\\
  -- Expand (a + b) squared using the binomial identity\\
\par\textbf{Compiler message excerpt:}\par\ttfamily
line 1: unknown module prefix 'algebra'\\
\\
No directory 'algebra' or file 'algebra.olean' in the search path entries:\\
/Users/sadiapeerzada/NLP-to-Formal-Language-Proof/lean\_project/.lake/packages/Cli/.lake/build/lib/lean\\
/Users/sadiapeerzada/NLP-to-Formal-Language-Proof/lean\_project/.lake/packages/batteries/.lake/build/lib/lean\\
/Users/sadiapeerzada/NLP-to-Formal-Language-Proof/lean\_project/.lake/packages/Qq/.lake/build/lib/lean\\
\par\textbf{Attempt 1 to attempt 2 candidate diff:}\par\ttfamily
--- attempt\_1\\
+++ attempt\_2\\
@@ -1,4 +1,3 @@\\
-import algebra.group\_power\\
 import algebra.ring\\
 \\
 theorem add\_sq\_nat (a b : ℕ) : (a + b) \textasciicircum{} 2 = a \textasciicircum{} 2 + 2 * a * b + b \textasciicircum{} 2 := by\\
\par\textbf{Saved verifier outcomes:} attempt 1 success=false; attempt 2 success=false. Exit codes are not recorded in this legacy trace.
\end{quote}
\paragraph{arith\_004 (manual trace; attempts 1 and 2 both rejected).}
\begin{quote}\small
\textbf{Initial candidate excerpt:}\par\ttfamily
theorem sq\_diff\_plus\_2ab (a b : ℕ) (h : b ≤ a) : (a - b) \textasciicircum{} 2 + 2 * a * b = a \textasciicircum{} 2 + b \textasciicircum{} 2 := by\\
  have h1 : (a - b) \textasciicircum{} 2 = a \textasciicircum{} 2 - 2 * a * b + b \textasciicircum{} 2 := by\\
    rw [Nat.sub\_def, Nat.mul\_add, Nat.sq\_diff]\\
  have h2 : a \textasciicircum{} 2 - 2 * a * b + b \textasciicircum{} 2 + 2 * a * b = a \textasciicircum{} 2 + b \textasciicircum{} 2 := by\\
    ring\\
\par\textbf{Compiler message excerpt:}\par\ttfamily
line 5: Unknown constant `Nat.sub\_def`\\
\par\textbf{Attempt 1 to attempt 2 candidate diff:}\par\ttfamily
--- attempt\_1\\
+++ attempt\_2\\
@@ -1,6 +1,6 @@\\
 theorem sq\_diff\_plus\_2ab (a b : ℕ) (h : b ≤ a) : (a - b) \textasciicircum{} 2 + 2 * a * b = a \textasciicircum{} 2 + b \textasciicircum{} 2 := by\\
   have h1 : (a - b) \textasciicircum{} 2 = a \textasciicircum{} 2 - 2 * a * b + b \textasciicircum{} 2 := by\\
-    rw [Nat.sub\_def, Nat.mul\_add, Nat.sq\_diff]\\
+    rw [Nat.sq\_diff]\\
   have h2 : a \textasciicircum{} 2 - 2 * a * b + b \textasciicircum{} 2 + 2 * a * b = a \textasciicircum{} 2 + b \textasciicircum{} 2 := by\\
\par\textbf{Saved verifier outcomes:} attempt 1 success=false; attempt 2 success=false. Exit codes are not recorded in this legacy trace.
\end{quote}
\paragraph{arith\_003 (manual trace; attempts 1 and 2 both rejected).}
\begin{quote}\small
\textbf{Initial candidate excerpt:}\par\ttfamily
theorem distributive\_nat (a b c : ℕ) : a * (b + c) = a * b + a * c :=\\
begin\\
  -- Step 1: Expand the product using the distributive property of multiplication over addition\\
  rw mul\_add,\\
  -- Step 2: No additional steps needed, the equation is already in the desired form\\
\par\textbf{Compiler message excerpt:}\par\ttfamily
line 4: Unknown identifier `begin`\\
\par\textbf{Attempt 1 to attempt 2 candidate diff:}\par\ttfamily
--- attempt\_1\\
+++ attempt\_2\\
@@ -1,12 +1,2 @@\\
-theorem distributive\_nat (a b c : ℕ) : a * (b + c) = a * b + a * c :=\\
-begin\\
-  -- Step 1: Expand the product using the distributive property of multiplication over addition\\
-  rw mul\_add,\\
-  -- Step 2: No additional steps needed, the equation is already in the desired form\\
\par\textbf{Saved verifier outcomes:} attempt 1 success=false; attempt 2 success=false. Exit codes are not recorded in this legacy trace.
\end{quote}
\paragraph{ineq\_001 (manual trace; attempts 1 and 2 both rejected).}
\begin{quote}\small
\textbf{Initial candidate excerpt:}\par\ttfamily
theorem at\_most\_n\_sq (n : ℕ) : n ≤ n \textasciicircum{} 2 := by\\
  -- Recall the definition of the square of a natural number\\
  have h1 : n \textasciicircum{} 2 = n * n := by\\
    -- Use the commutative property of multiplication to rewrite n squared as n * n\\
    rw [mul\_comm]\\
\par\textbf{Compiler message excerpt:}\par\ttfamily
line 5: unsolved goals\\
n : ℕ\\
⊢ n \textasciicircum{} 2 = n * n\\
\par\textbf{Attempt 1 to attempt 2 candidate diff:}\par\ttfamily
--- attempt\_1\\
+++ attempt\_2\\
@@ -7,7 +7,7 @@\\
   -- Expand n * n using the distributive property of multiplication over addition\\
   have h2 : n * n = (1 + n) * n := by\\
     -- Since n is a natural number, it can be written as the sum of 1 and itself\\
-    rw [add\_comm, add\_self\_add\_right]\\
+    rw [add\_comm, add\_one]\\
\par\textbf{Saved verifier outcomes:} attempt 1 success=false; attempt 2 success=false. Exit codes are not recorded in this legacy trace.
\end{quote}
\paragraph{mod\_002 (manual trace; attempts 1 and 2 both rejected).}
\begin{quote}\small
\textbf{Initial candidate excerpt:}\par\ttfamily
theorem sq\_mod\_four (n : ℤ) : n \textasciicircum{} 2 \% 4 = 0 ∨ n \textasciicircum{} 2 \% 4 = 1 :=\\
begin\\
  cases (Int.even\_or\_odd n) with k hk,\\
  \{\\
    -- n is even\\
\par\textbf{Compiler message excerpt:}\par\ttfamily
line 4: Unknown identifier `begin`\\
\par\textbf{Attempt 1 to attempt 2 candidate diff:}\par\ttfamily
--- attempt\_1\\
+++ attempt\_2\\
@@ -1,22 +1,17 @@\\
+-- The previous attempt failed due to incorrect syntax for cases and begin/end blocks.\\
 theorem sq\_mod\_four (n : ℤ) : n \textasciicircum{} 2 \% 4 = 0 ∨ n \textasciicircum{} 2 \% 4 = 1 :=\\
-begin\\
-  cases (Int.even\_or\_odd n) with k hk,\\
-  \{\\
\par\textbf{Saved verifier outcomes:} attempt 1 success=false; attempt 2 success=false. Exit codes are not recorded in this legacy trace.
\end{quote}
